\section{Introducción}

El laboratorio de análisis industrial de BRINSA en Cajicá (Cundinamarca) requiere un sistema de ventilación por presurización positiva cuya función es doble: mantener el recinto en presión positiva respecto al exterior y excluir insectos, objetos extraños y polvo del ambiente del laboratorio. El sistema impulsa 3 840 m$^3$/h de aire filtrado, equivalentes a 12 renovaciones/h sobre un volumen efectivo de 320 m$^3$, y mantiene un diferencial de presión con set-point de +25 Pa y mínimo admisible de +12.5 Pa. No se trata de un sistema de biocontención; por tanto, se descarta expresamente la filtración HEPA y la etapa definitiva de filtración es MERV 13-14.

La selección de los equipos del sistema presenta dos condicionantes que la distinguen de una especificación convencional. El primero es la altitud del sitio (2 558 msnm, presión atmosférica 74.1 kPa, densidad del aire 0.88 kg/m$^3$), que obliga a corregir por densidad todos los valores de presión y potencia publicados por los fabricantes a densidad estándar. El segundo es el ambiente exterior de la planta de hipoclorito de calcio, altamente corrosivo, que restringe los materiales admisibles para el ventilador, las rejillas, el damper de alivio, la ferretería y los sellantes.

El presente informe documenta la investigación de mercado y de fuentes técnicas ejecutada para seleccionar y justificar cada componente del sistema. El documento se estructura de la siguiente manera. La Sección 2 presenta los objetivos y la Sección 3 el alcance. La Sección 4 define las bases de diseño: condiciones ambientales del sitio, corrección de densidad y corrosividad del ambiente. La Sección 5 describe la metodología de investigación y los criterios de aceptación de fuentes. La Sección 6 presenta la descripción funcional del sistema. Las Secciones 7 y 8 desarrollan el análisis comparativo por componente: ventilador y filtración en la primera; rejillas de exfiltración, damper de alivio, instrumentación y accesorios en la segunda. La Sección 9 compila el marco normativo aplicable. Finalmente, las Secciones 10 y 11 presentan las conclusiones y la selección recomendada.
